\section{Source review and the precise new contribution}
\label{sec:sources}

\subsection{What was taken from the three requested sources}

The upstream repository was inspected at commit
\texttt{72fead102f663d5c7a407964a13dc0cf9d977207}.  Its
\texttt{grh-compact/src/Statements.lean} imports Mathlib and defines the
zero-location statements, while \texttt{src/Main.lean} proves those
statements and Mathlib's \texttt{RiemannHypothesis}.  The scope is the
Riemann zeta function and \emph{Dirichlet} $L$-functions, including all
positive moduli and imprimitive characters.  It does not state GRH for
arbitrary Artin or automorphic $L$-functions~\cite{upstream}.

The repository reports successful compilation, kernel replay, and
Comparator checks, with the five main theorems depending only on
\texttt{propext}, \texttt{Classical.choice}, and \texttt{Quot.sound}.
The readable and compact editions contain roughly 300,000 and 76,000
lines respectively.  The pinned source and selected central files were
read for this report; the complete build was not replayed here.
Andrew Sutherland's verification page records independent checks of this
same shorter release~\cite{sutherland}.  The requested Wikipedia page
reports the November 2026 proof and its verification~\cite{wikipedia};
the exact formal statements, rather than the encyclopedia's broad
terminology, determine the scope used below.

The source's central new argument is an \emph{arithmetic covolume
comparison}, not the spectral construction in this report.  In the
compact edition, \texttt{New/Endgame/ShortAndEnough.lean} supplies many
small integral forms and an almost-spanning dimension bound;
\texttt{LongForms.lean} supplies a lower bound for the same integral
covolume using heights and the product formula.  In the notation of
\texttt{New/Main.lean}, cancellation of the main height leads to
\[
 mR\le (\log b)^2\,\mathrm{budget},\qquad m\asymp\sqrt x,
\]
where $R$ is a smoothed prime-power remainder.  The compact
\texttt{OneSidedCriterion.lean} then applies a one-sided Landau argument;
functional equations complete the zero-location conclusion.
The useful conceptual pattern is to retain positivity and a common
determinant while comparing two descriptions of it.  Our argument
uses that pattern in a new analytic construction after importing the
proved zero-location statement.

The ProveIt tree was inspected at commit
\texttt{041d8032541746f2ac2e686fbe437f8092771670}.  Its README distinguishes
61 imported older research packets from independent checking and states
that their Lean excerpts are not part of its root Lake package
\cite{proveit}.  The review here read the index, relevant merged-article
sections, the graph packet, the robust-spectrum packet, and descriptions
and targeted source searches in seven later incoming packets.  It is
not a proof audit of every report in the collection.  Retrieved archive
bytes were checked against their pinned Git blob identifiers, and a
source manifest is included with this report.

\subsection{The immediate predecessor}

The packet \texttt{graph\_coupled\_zeros\_2026-11-26.zip} already proves
the finite polynomial identity
\[
 P_G(z)=\det\!\left(\operatorname{diag}(p_i(z))L_G+
                             \operatorname{diag}(p_i'(z))\right)
\]
has real zeros, with a compression interpretation and an exact total
displacement expressed through effective resistances.  It also proves
real-zero preservation, collision multiplicities, and bounded zero-count
discrepancy for entire inputs~\cite{priorgraph}.  Those facts are prior
work; the essential half-plane argument is reproduced below for a
self-contained proof of the extensions.

Its final section explicitly asks whether one can regularize the
displacement of \emph{all} completed Dirichlet zeros and obtain an
infinite analogue of the electrical formula.  Theorems
\ref{ut:exact}--\ref{ut:expansion} answer the convergence and asymptotic
part for the original construction.  Theorem~\ref{thm:electrical-infinite}
gives an electrical counterpart for a separately defined squared-ordinate
construction.  In particular, an unweighted infinite sum is generally
\emph{divergent}; its precise growth and convergent transforms are the
correct objects.

\subsection{Novelty compared with primary literature}

The mathematical components have clear precedents.  Matrix meromorphic
Herglotz theory already provides higher-order interlacing
\cite{reiffenstein}.  Spectral-shift logarithms and trace identities are
classical~\cite{bgn}.  Herglotz boundary distributions and Cauchy laws are
treated by Aizenman and Warzel~\cite{aizenmanwarzel}.  Sodin proves a weak
Cauchy limit for the normalized logarithmic derivative under RH
\cite[Section 3.4.5, equation (55)]{sodin}; Najnudel and Nikeghbali develop
that theme further~\cite{najnudel}.  Differencing and derivative
deformations of completed $L$-functions also predate this work
\cite{lagarias,farmergonek}.

The contributions developed and proved here are the following:
\begin{enumerate}
\item An arithmetic local limit for the logarithmic derivative valid for
\emph{arbitrarily shrinking moving value intervals}, including relative
asymptotics for exponentially small widths
(Theorem~\ref{thm:local-cauchy}).
\item A universal local transport limit for positive Hermitian
perturbations of completed Dirichlet determinants, on every fixed
macroscopic real window (Theorem~\ref{ut:main}).
\item An all-orders Poisson-smoothed expansion whose coefficients are
explicit finite matrix traces, separating the first conductor and graph
corrections (Theorem~\ref{ut:expansion}).
\item A squared-ordinate transport law with exact coefficient given by
the decrease in the graph's Kirchhoff index
(Theorem~\ref{thm:electrical-infinite}).
\end{enumerate}
No identical statement was located in the inspected ProveIt material or
the targeted primary-literature search.  This is a documented novelty
comparison, not certification of worldwide publication priority.  The
proofs below establish the claims as mathematical statements independently
of that priority question.

\paragraph{Dependence on RH/GRH.}
The given zero-location theorem makes the relevant logarithmic derivatives
have a strict half-plane sign.  The functional equation and gamma-factor
asymptotics then determine the constants.  No simplicity of zeros,
linear independence of ordinates, pair-correlation conjecture, or
identification of the new determinants with Euler products is assumed.
The proofs are written out in ordinary mathematics.  They have not been
formalized in Lean in this packet.
